\documentclass[a4paper,10pt]{report}\input{../0.Préambule/Préambule_cours_prof}\begin{document}

\newpage
\section*{ALGORITHMIQUE ET PROGRAMMATION - Travaux pratiques 1}
\addcontentsline{toc}{section}{ALGORITHMIQUE ET PROGRAMMATION - Travaux pratiques 1}

\begin{dinglist}{111}
\item Recopie le script de ton cours qui te permettra de débuter tous les tracés, à toi de t'adapter si besoin.
\end{dinglist}

\subsection*{Mes premières figures}
\addcontentsline{toc}{section}{Mes premières figures}

\subsubsection{Le carré}

\begin{dinglist}{111}
\item Écris un programme qui permet de tracer un carré de coté $100$ pixels avec la brique \textbf{répéter} : 
\item Dessine la figure obtenu et indique par un croix, le point de départ et d'arrivée du lutin :

\vspace{-2cm}
\begin{flushright}
\begin{scratch}[scale=0.8]
\blockrepeat{répéter \ovalnum{} fois}
{
\blockspace
}
\end{scratch}
\end{flushright}

\renewcommand{\arraystretch}{2}
\begin{tabular}{m{4cm}m{2.4cm}m{1cm}m{4cm}m{2.4cm}}
$\bullet$ quand il tourne à gauche : &
\grandscarreaux{\linewidth}{1.6cm}&&
$\bullet$ quand il tourne à droite : &
\noindent\grandscarreaux{\linewidth}{1.6cm}\\
\end{tabular}

\ding{226} \textit{Garde une trace de ce programme sur ton écran pour la vérification}
\end{dinglist}

\subsubsection{Le rectangle}

\begin{dinglist}{111}

\item Toujours avec le bloc \textbf{répéter}, écris un autre programme qui permet de tracer un rectangle (horizontal) de longueur $100$ et de largeur $50$.

\ding{226} \textit{Garde une trace de ce programme sur ton écran pour la vérification}

\vspace{3mm}
\item Modifie le script pour obtenir un rectangle (vertical) de même dimension que le précédent.

\ding{226} \textit{Garde une trace de ce programme sur ton écran pour la vérification}
\end{dinglist}

\subsubsection{Le triangle équilatéral}

\begin{dinglist}{111}
\item Donne la mesure des angles d'un triangle équilatéral.

\medskip
\noindent\grandscarreaux{\linewidth}{0.8cm}

\item L'angle a indiquer dans Scratch est égal à la mesure de l'angle plat moins l'angle à tracer. De combien doit tourner le lutin ?

\medskip
\noindent\grandscarreaux{\linewidth}{0.8cm}

\item Écris deux programmes qui permettent de tracer ces deux triangles équilatéraux de coté $100$.

\ding{226} \textit{Garde une trace de ce programme sur ton écran pour la vérification}

\ding{226} \textit{Appelle le professeur pour la vérification}
\end{dinglist}

\subsection*{Autres figures régulières}
\addcontentsline{toc}{section}{Autres figures régulières}

\begin{dinglist}{111}
\item Écris un programme permettant d'obtenir un pentagone régulier.
\item Écris un programme permettant d'obtenir un hexagone régulier.
\item Pourrais-tu trouver une relation entre l'angle à indiquer dans Scratch et le nombre de coté du polygone ?

\medskip
\noindent\grandscarreaux{\linewidth}{1.6cm}
\item Écris un programme permettant d'obtenir un polygone à plusieurs cotés, le nombre de cotés étant demandé à l'utilisateur.

\ding{226} \textit{Appelle le professeur pour la vérification}

\item Écris un programme permettant d'obtenir un cercle.

\item Réalise un rosace.
\end{dinglist}

\end{document}